\documentclass[10pt,a4paper,twocolumn]{article}
\usepackage[a4paper,margin=17mm,columnsep=6mm]{geometry}
\usepackage[T1]{fontenc}
\usepackage{lmodern}
\usepackage{microtype}
\usepackage{amsmath}
\usepackage{booktabs}
\usepackage{siunitx}
\usepackage{pgfplots}
\usepackage{authblk}
\usepackage[hidelinks]{hyperref}
\pgfplotsset{compat=1.18}
\sisetup{group-separator={\,},group-minimum-digits=4}
\newcommand{\code}[1]{\texttt{#1}}
\newcommand{\CxErr}{\num{2.0e-7}}
\newcommand{\CxErrD}{\num{5.6e-16}}
\newcommand{\DoubleCxLarge}{0.88--0.92}
\newcommand{\DoubleCxSmall}{0.88--1.27}
\newcommand{\DoubleCxbcLarge}{0.87--0.93}
\newcommand{\DoubleCxbcSmall}{0.85--1.11}
\newcommand{\DoubleCxxLarge}{0.97--1.29}
\newcommand{\DoubleCxxSmall}{1.09--1.17}
\newcommand{\DoubleErr}{\num{6.0e-16}}
\newcommand{\DoubleErrGr}{\num{8.4e-16}}
\newcommand{\DoubleTrunkVsGr}{1.4--2.0}
\newcommand{\DoubleVsFftwRange}{0.70--1.20$\times$}
\newcommand{\DoubleVsGrRange}{1.3--2.0$\times$}
\newcommand{\FallbackSlowdown}{3.6--11.3$\times$}
\newcommand{\FallbackSlowdownGr}{1.2--2.0$\times$}
\newcommand{\FloatCxLarge}{0.88--0.95}
\newcommand{\FloatCxSmall}{1.13--1.48}
\newcommand{\FloatCxbcLarge}{0.83--0.93}
\newcommand{\FloatCxbcSmall}{0.99--1.18}
\newcommand{\FloatCxxLarge}{0.98--1.18}
\newcommand{\FloatCxxSmall}{1.05--1.25}
\newcommand{\FloatErr}{\num{2.1e-7}}
\newcommand{\FloatErrGr}{\num{4.8e-7}}
\newcommand{\FloatTrunkVsGr}{1.6--2.9}
\newcommand{\FloatVsFftwRange}{0.86--2.04$\times$}
\newcommand{\FloatVsGrRange}{1.6--3.4$\times$}
\newcommand{\GrSpectrumComplex}{1.06$\times$}
\newcommand{\GrSpectrumReal}{1.00$\times$}
\newcommand{\GrStreamDouble}{1.3--1.8$\times$}
\newcommand{\GrStreamFloat}{1.7--2.5$\times$}
\newcommand{\HostCPU}{Intel Xeon Processor @ 2.10GHz}
\newcommand{\HostCompiler}{g++-14 14.2.0}
\newcommand{\HostCores}{4 vCPU}
\newcommand{\HostKernel}{6.18.44-fc-v114}
\newcommand{\Samples}{120}
\newcommand{\TrunkCommit}{575b8ec6ec}
\newcommand{\TrunkVersion}{17.0.0 20261010}
\newcommand{\WidthGain}{1.4--3.0}
\newcommand{\WidthSixteenVsGr}{0.9--1.8$\times$}
\newcommand{\WorstMAD}{10.2\,\%}
% measured values, generated by mkdata.py from ../results

\title{\vspace{-8mm}\bfseries A portable \texttt{std::simd} FFT\\ 2--4$\times$ faster than GNU~Radio~4's}
\author[1]{Dennis Klein}
\author[2]{Claude}
\affil[2]{Anthropic}
\date{\vspace{-4mm}\small October 2026}

\begin{document}
\maketitle

\begin{abstract}
\noindent GNU Radio~4 computes FFTs with \code{SimdFFT}, a \code{std::simd} port of PFFFT whose vectors are fixed
at four lanes. We present \code{fastfft}, a single-header radix-8/16 Stockham FFT written against
\code{native\_simd} without ISA intrinsics. On one AVX-512 core it is \FloatVsGrRange{} faster than \code{SimdFFT} in
single precision (N\,=\,64\ldots$2^{20}$) and \DoubleVsGrRange{} in double precision, and within
\FloatVsFftwRange{} of FFTW~3.3.10 with AVX-512 codelets. Most of the gain comes from using the full vector
width; the rest from fewer passes over memory. Getting there with \code{std::simd} took three workarounds that we
report. A patch puts the kernel behind GNU~Radio~4's FFT API; the code and a harness to repeat all measurements on
any Linux host are public.
\end{abstract}

\section{Introduction}
GNU Radio~4 (gr4)~\cite{gr4} is a header-only C++23 rewrite of the GNU Radio signal-processing framework. Its FFT
blocks call \code{gr::algorithm::FFT}, which dispatches to \code{SimdFFT}, a port of PFFFT~\cite{pffft} to
\code{std::experimental::simd}~\cite{ts} via vir-simd~\cite{vir}. \code{SimdFFT} uses
\code{simd<T,\,4>}: four floats or four doubles, whatever the hardware offers. On an AVX-512 machine this leaves
three quarters of each register unused in single precision.

We asked how fast a power-of-two complex FFT can be when written only against \code{native\_simd}, so that the
width follows the compilation target, and whether it can match a tuned library. The answer, \code{fastfft}, is one
\SI{400}{line} header. This note describes its structure (\S\ref{sec:alg}), what \code{std::simd} could not express
(\S\ref{sec:simd}), the measurement method (\S\ref{sec:method}) and the results (\S\ref{sec:results}).

\section{Algorithm}\label{sec:alg}
\code{fastfft} is a Stockham autosort FFT~\cite{stockham,vanloan}: every pass reads one buffer and writes the other
in natural order, so no bit-reversal pass is needed. Data are interleaved complex numbers; a vector of $W$ scalars
holds $C=W/2$ complex values (8 floats or 4 doubles per \SI{512}{bit} register).
\begin{enumerate}\setlength\itemsep{0pt}
\item \emph{First pass, radix 8, vectorised over the butterfly index.} With stride 1 the eight inputs of a butterfly
  are $N/8$ apart, so $C$ consecutive butterflies form one vector per input. The eight results of each butterfly
  must be stored contiguously; an in-register $8\times C$ complex transpose (three rounds of two-input zips) does
  this before the store.
\item \emph{Middle passes, radix 16, vectorised over the stride.} Once the stride $s\ge C$, the $C$ lanes of a vector
  are $C$ independent sub-transforms that share a twiddle factor, which is broadcast. A radix-16 butterfly is done as
  $4\times4$ in registers. Each pass reads and writes the data once, so a $2^{20}$ transform takes six passes
  instead of ten for radix 4.
\item \emph{Last pass,} radix 2, 4 or 8, whatever remains.
\end{enumerate}
Complex multiplication by a twiddle $w$ uses tables stored as $(w_r,w_r)$ and $(-w_i,w_i)$ per lane pair, so
$z\cdot w = z\odot\mathbf{w_r} + \mathrm{swap}(z)\odot\mathbf{w_i'}$ costs one multiply, one fused multiply-add and
one in-lane swap. Twiddles are precomputed in double precision. Sizes below 64 use a scalar radix-2 path.

\section{What \texttt{std::simd} could not express}\label{sec:simd}
The kernel uses \code{std::experimental::simd} as shipped with libstdc++ and contains no intrinsics, but three
points needed care.
\textbf{(i) Two-input shuffles.} The transpose in pass~1 needs zip operations that take lanes from two vectors.
\code{std::simd} has none; \code{concat} followed by a permute compiled to element-wise moves that cost about three
times the whole pass. We use the compiler's generic vector extension (\code{\_\_builtin\_shufflevector}), which
GCC and clang lower to single \code{vpermt2ps}-class instructions on every target. The re/im swap is written the
same way.
\textbf{(ii) FMA.} libstdc++'s \code{fma(simd,simd,simd)} is not inlined; we write \code{a*b+c} and compile with
\code{-ffp-contract=fast} (GCC's default in GNU mode).
\textbf{(iii) Backends.} vir-simd forwards to libstdc++'s implementation only when it is available; with clang and
libc++ it falls back to array-based emulation. There \code{fastfft} is \FallbackSlowdown{} slower, and gr4's own
\code{SimdFFT} is about 2$\times$ slower than with GCC (both measured, appendix). A gr4 built with
clang\,+\,libc++ therefore loses most of its SIMD performance without notice. Our header defines
\code{FASTFFT\_NATIVE\_SIMD} so that callers can detect this case.

\section{Method}\label{sec:method}
\code{fftbench} times \code{fastfft}, \code{SimdFFT} (gr4 rev.~\code{af62235}) and FFTW~3.3.10~\cite{fftw}
(built from source with the host's SIMD codelets, \code{FFTW\_MEASURE}) on the same input, uniform random in
$[-1,1)^2$. A \emph{sample} transforms $\max(N,2^{16})$ consecutive values as frames of $N$, the way a streaming gr4
block sees them, repeated until it lasts at least \SI{5}{ms}; one warm-up sample is discarded. Before timing, every
output is checked against a long-double reference DFT. Each process runs 40 rounds; in every round each
implementation is sampled once in a freshly randomised order, so drifts affect all of them alike and same-round
samples can be paired. The process is run three times and pinned to one core. We report the median of the 120
samples per cell with a 95\,\% percentile-bootstrap confidence interval, and speedups as the median of paired
per-round ratios, also with a bootstrap CI~\cite{efron,hoefler}. All values, raw samples and host data are in the
repository; one script (\code{bench/run.sh}) repeats the whole procedure.

Host: \HostCPU{}, \HostCores{}, Linux~\HostKernel{}, \HostCompiler{}, \code{-O3 -march=native}.
The machine is a virtual machine; the largest relative median absolute deviation over all cells was \WorstMAD{}.

\section{Results}\label{sec:results}
\begin{figure}[t]
\centering
\begin{tikzpicture}
\begin{axis}[width=\columnwidth,height=52mm,xmode=log,log basis x=2,ymode=log,log basis y=2,
  xlabel={transform size $N$},ylabel={speedup of \code{fastfft}},ytick={0.5,1,2,4},yticklabels={0.5,1,2,4},
  xtick={64,256,1024,4096,16384,65536,262144,1048576},
  xticklabels={$2^6$,$2^8$,$2^{10}$,$2^{12}$,$2^{14}$,$2^{16}$,$2^{18}$,$2^{20}$},
  ymin=0.45,ymax=6,grid=major,grid style={gray!25},
  legend style={font=\scriptsize,at={(0.02,0.02)},anchor=south west,legend columns=2,draw=none,fill=white,fill opacity=.8,text opacity=1},
  tick label style={font=\scriptsize},label style={font=\small}]
\addplot+[mark=*,mark size=1.3pt,error bars/.cd,y dir=both,y explicit] table[x=n,y=s,y error plus=ep,y error minus=em]{data/float-gr4.dat};
\addplot+[mark=square*,mark size=1.3pt,error bars/.cd,y dir=both,y explicit] table[x=n,y=s,y error plus=ep,y error minus=em]{data/double-gr4.dat};
\addplot+[mark=o,mark size=1.5pt,dashed,error bars/.cd,y dir=both,y explicit] table[x=n,y=s,y error plus=ep,y error minus=em]{data/float-fftw.dat};
\addplot+[mark=square,mark size=1.5pt,dashed,error bars/.cd,y dir=both,y explicit] table[x=n,y=s,y error plus=ep,y error minus=em]{data/double-fftw.dat};
\legend{vs.\ SimdFFT, float,vs.\ SimdFFT, double,vs.\ FFTW, float,vs.\ FFTW, double}
\end{axis}
\end{tikzpicture}
\vspace{-2mm}
\caption{Speedup of \code{fastfft} (median of paired ratios, 95\,\% CI as error bars; mostly smaller than the
markers). Above 1 means \code{fastfft} is faster. Absolute times in Tables~\ref{tab:float} and~\ref{tab:double}.}
\label{fig:speedup}
\end{figure}

\paragraph{Kernel.} Figure~\ref{fig:speedup} summarises the comparison. In single precision \code{fastfft} is
\FloatVsGrRange{} faster than \code{SimdFFT} over the whole range; the advantage shrinks above $2^{16}$, where both are
limited by the memory hierarchy. Against FFTW it ranges from \FloatVsFftwRange{}: ahead at small sizes and at
$2^{19}$--$2^{20}$, behind around $2^{12}$ and $2^{18}$. In double precision the gain over \code{SimdFFT} is
\DoubleVsGrRange{}, smaller because four doubles already fill half of an AVX-512 register; FFTW is
\DoubleVsFftwRange{}. Accuracy is unchanged: the largest error relative to $\max|X|$ was \FloatErr{} (float) and
\DoubleErr{} (double), against \FloatErrGr{} and \DoubleErrGr{} for \code{SimdFFT}.

\paragraph{Vector width.} Forcing \code{fastfft} to narrower vectors isolates the width effect
(Table~\ref{tab:width}). At \SI{128}{bit}, the width \code{SimdFFT} uses for floats, \code{fastfft} is still
\WidthSixteenVsGr{} faster than \code{SimdFFT}, which is the share due to the algorithm; going to \SI{512}{bit} adds a further
factor of \WidthGain{}.

\paragraph{In gr4.} With the patch, \code{gr::algorithm::FFT} uses \code{fastfft} for complex power-of-two sizes
$\ge 64$ when the \code{simd} backend is native. The spectrum block (\code{FFT<T,DataSet<P>>}) also computed into a
freshly allocated vector per frame and made two extra passes for the fftshift; the patch removes both. gr4's own
\code{bm\_fft\_block} (N\,=\,1024, unchanged) improves by \GrBlockComplex{} for complex input; real input still uses
\code{SimdFFT}'s real transform and gains \GrBlockReal{} from the buffer changes alone. Most of the remaining time
per frame is spent outside the FFT, in \code{atan2}, \code{hypot} and building the \code{DataSet}. gr4's unit tests
for the FFT blocks and algorithm pass unchanged.

\section{Discussion}
Using the native width and fewer, larger radix passes is enough to make a portable FFT competitive with a
hand-tuned library on the measured host; nothing in the kernel is specific to x86. What \code{std::simd} lacks is a
two-input permute; with that added (or the C++26 \code{std::simd} providing one), the kernel would be pure standard
C++. The gap to FFTW at some mid sizes indicates where planning (choosing radices per size) would help.

These results come from one virtual AVX-512 core. Wider claims need more platforms---AVX2-only x86, ARM with NEON
or SVE---and the harness was written so that collecting them takes one command. The kernel covers what gr4's FFT
block needs most (forward, complex, power of two); real-input and non-power-of-two transforms still use
\code{SimdFFT}.

\paragraph{Code.} \url{https://github.com/dennisklein/fft} (header, harness, gr4 patch, raw data).

\paragraph{Contributions.} D.K.\ posed the question, steered the work and reviewed the results. Claude, an AI model
by Anthropic, wrote the code, ran the measurements and drafted the text.

\begin{thebibliography}{9}\setlength\itemsep{0pt}\small
\bibitem{gr4} GNU Radio 4.0, \url{https://github.com/fair-acc/gnuradio4}, rev.\ \code{af62235}.
\bibitem{pffft} J.~Pommier, PFFFT: a pretty fast FFT, \url{https://bitbucket.org/jpommier/pffft}.
\bibitem{ts} ISO/IEC TS 19570:2018, Parallelism TS v2 (\code{std::experimental::simd}); P1928, \code{std::simd} for C++26.
\bibitem{vir} M.~Kretz, vir-simd, \url{https://github.com/mattkretz/vir-simd}, v0.4.4.
\bibitem{fftw} M.~Frigo, S.~G.~Johnson, The design and implementation of FFTW3, Proc.\ IEEE 93(2):216--231, 2005.
\bibitem{stockham} T.~G.~Stockham, High-speed convolution and correlation, AFIPS SJCC, 1966.
\bibitem{vanloan} C.~Van Loan, Computational Frameworks for the Fast Fourier Transform, SIAM, 1992.
\bibitem{efron} B.~Efron, R.~J.~Tibshirani, An Introduction to the Bootstrap, Chapman \& Hall, 1993.
\bibitem{hoefler} T.~Hoefler, R.~Belli, Scientific benchmarking of parallel computing systems, SC'15, 2015.
\end{thebibliography}

\clearpage
\onecolumn
\appendix
\section*{Appendix: measurement data}
\small
\begin{table}[h]\centering
\caption{Complex float: median ns per transform [95\,\% CI] and speedup of fastfft (median of paired ratios [95\,\% CI]).}\label{tab:float}
\begin{tabular}{rlllll}\toprule
$N$ & fastfft & gr4 SimdFFT & FFTW 3.3.10 & vs.\ SimdFFT & vs.\ FFTW\\\midrule
$2^{6}$ & 35.0 {\scriptsize[34.7, 35.2]} & 114 {\scriptsize[112, 116]} & 44.1 {\scriptsize[43.5, 44.6]} & 3.27 {\scriptsize[3.23, 3.30]} & 1.26 {\scriptsize[1.25, 1.27]}\\
$2^{7}$ & 69.0 {\scriptsize[68.5, 69.7]} & 237 {\scriptsize[235, 238]} & 76.7 {\scriptsize[76.3, 77.1]} & 3.42 {\scriptsize[3.40, 3.44]} & 1.11 {\scriptsize[1.10, 1.12]}\\
$2^{8}$ & 162 {\scriptsize[161, 164]} & 468 {\scriptsize[464, 471]} & 162 {\scriptsize[160, 163]} & 2.89 {\scriptsize[2.87, 2.92]} & 0.99 {\scriptsize[0.98, 1.00]}\\
$2^{9}$ & 302 {\scriptsize[298, 305]} & 1018 {\scriptsize[1010, 1031]} & 356 {\scriptsize[353, 362]} & 3.39 {\scriptsize[3.36, 3.41]} & 1.16 {\scriptsize[1.16, 1.18]}\\
$2^{10}$ & 655 {\scriptsize[649, 659]} & 2147 {\scriptsize[2134, 2159]} & 759 {\scriptsize[752, 764]} & 3.29 {\scriptsize[3.25, 3.32]} & 1.17 {\scriptsize[1.16, 1.17]}\\
$2^{11}$ & 1793 {\scriptsize[1784, 1807]} & 4942 {\scriptsize[4917, 4959]} & 1696 {\scriptsize[1682, 1708]} & 2.76 {\scriptsize[2.75, 2.77]} & 0.95 {\scriptsize[0.94, 0.96]}\\
$2^{12}$ & 4895 {\scriptsize[4866, 4922]} & 12.2\,k {\scriptsize[12.1\,k, 12.2\,k]} & 3854 {\scriptsize[3805, 3892]} & 2.50 {\scriptsize[2.48, 2.51]} & 0.79 {\scriptsize[0.78, 0.79]}\\
$2^{13}$ & 10.1\,k {\scriptsize[9997, 10.2\,k]} & 29.3\,k {\scriptsize[29.2\,k, 29.4\,k]} & 12.6\,k {\scriptsize[12.5\,k, 12.7\,k]} & 2.91 {\scriptsize[2.89, 2.93]} & 1.25 {\scriptsize[1.24, 1.26]}\\
$2^{14}$ & 20.5\,k {\scriptsize[20.3\,k, 20.6\,k]} & 59.1\,k {\scriptsize[58.8\,k, 59.5\,k]} & 28.1\,k {\scriptsize[27.9\,k, 28.2\,k]} & 2.90 {\scriptsize[2.88, 2.93]} & 1.36 {\scriptsize[1.34, 1.38]}\\
$2^{15}$ & 45.6\,k {\scriptsize[45.2\,k, 46.2\,k]} & 134.7\,k {\scriptsize[134.2\,k, 135.7\,k]} & 62.0\,k {\scriptsize[61.5\,k, 62.6\,k]} & 2.96 {\scriptsize[2.94, 2.98]} & 1.36 {\scriptsize[1.34, 1.38]}\\
$2^{16}$ & 140.9\,k {\scriptsize[140.4\,k, 141.7\,k]} & 271.7\,k {\scriptsize[269.7\,k, 273.6\,k]} & 139.1\,k {\scriptsize[138.4\,k, 139.9\,k]} & 1.93 {\scriptsize[1.89, 1.97]} & 0.98 {\scriptsize[0.97, 1.00]}\\
$2^{17}$ & 418.3\,k {\scriptsize[410.7\,k, 423.7\,k]} & 726.4\,k {\scriptsize[719.1\,k, 739.8\,k]} & 370.4\,k {\scriptsize[365.2\,k, 375.2\,k]} & 1.77 {\scriptsize[1.74, 1.79]} & 0.90 {\scriptsize[0.88, 0.92]}\\
$2^{18}$ & 1.008\,M {\scriptsize[1.001\,M, 1.015\,M]} & 1.808\,M {\scriptsize[1.792\,M, 1.832\,M]} & 881.0\,k {\scriptsize[873.5\,k, 890.8\,k]} & 1.80 {\scriptsize[1.78, 1.81]} & 0.88 {\scriptsize[0.87, 0.90]}\\
$2^{19}$ & 2.068\,M {\scriptsize[2.059\,M, 2.082\,M]} & 4.355\,M {\scriptsize[4.319\,M, 4.402\,M]} & 4.238\,M {\scriptsize[4.187\,M, 4.304\,M]} & 2.11 {\scriptsize[2.09, 2.13]} & 2.04 {\scriptsize[2.01, 2.06]}\\
$2^{20}$ & 5.524\,M {\scriptsize[5.439\,M, 5.616\,M]} & 9.040\,M {\scriptsize[8.948\,M, 9.148\,M]} & 8.338\,M {\scriptsize[8.246\,M, 8.428\,M]} & 1.64 {\scriptsize[1.63, 1.66]} & 1.53 {\scriptsize[1.51, 1.54]}\\
\bottomrule\end{tabular}\end{table}

\begin{table}[h]\centering
\caption{Complex double: median ns per transform [95\,\% CI] and speedup of fastfft (median of paired ratios [95\,\% CI]).}\label{tab:double}
\begin{tabular}{rlllll}\toprule
$N$ & fastfft & gr4 SimdFFT & FFTW 3.3.10 & vs.\ SimdFFT & vs.\ FFTW\\\midrule
$2^{6}$ & 72.8 {\scriptsize[71.3, 75.2]} & 130 {\scriptsize[128, 132]} & 84.5 {\scriptsize[84.0, 84.9]} & 1.79 {\scriptsize[1.78, 1.80]} & 1.13 {\scriptsize[1.12, 1.17]}\\
$2^{7}$ & 160 {\scriptsize[159, 162]} & 286 {\scriptsize[284, 288]} & 193 {\scriptsize[191, 195]} & 1.80 {\scriptsize[1.78, 1.82]} & 1.20 {\scriptsize[1.19, 1.21]}\\
$2^{8}$ & 383 {\scriptsize[380, 386]} & 596 {\scriptsize[589, 604]} & 396 {\scriptsize[392, 400]} & 1.57 {\scriptsize[1.55, 1.57]} & 1.03 {\scriptsize[1.02, 1.03]}\\
$2^{9}$ & 821 {\scriptsize[811, 829]} & 1367 {\scriptsize[1352, 1376]} & 866 {\scriptsize[854, 890]} & 1.67 {\scriptsize[1.65, 1.68]} & 1.05 {\scriptsize[1.04, 1.09]}\\
$2^{10}$ & 2155 {\scriptsize[2146, 2166]} & 3218 {\scriptsize[3202, 3236]} & 1857 {\scriptsize[1824, 1944]} & 1.49 {\scriptsize[1.48, 1.51]} & 0.85 {\scriptsize[0.84, 0.88]}\\
$2^{11}$ & 4952 {\scriptsize[4927, 4985]} & 10.1\,k {\scriptsize[10.0\,k, 10.2\,k]} & 4406 {\scriptsize[4314, 4482]} & 2.04 {\scriptsize[2.02, 2.05]} & 0.87 {\scriptsize[0.86, 0.89]}\\
$2^{12}$ & 12.2\,k {\scriptsize[12.1\,k, 12.3\,k]} & 18.6\,k {\scriptsize[18.4\,k, 18.7\,k]} & 11.4\,k {\scriptsize[11.3\,k, 11.5\,k]} & 1.52 {\scriptsize[1.51, 1.53]} & 0.93 {\scriptsize[0.92, 0.95]}\\
$2^{13}$ & 24.2\,k {\scriptsize[24.1\,k, 24.3\,k]} & 40.1\,k {\scriptsize[39.9\,k, 40.3\,k]} & 23.9\,k {\scriptsize[23.6\,k, 24.3\,k]} & 1.66 {\scriptsize[1.65, 1.67]} & 1.00 {\scriptsize[0.99, 1.00]}\\
$2^{14}$ & 51.5\,k {\scriptsize[51.2\,k, 51.8\,k]} & 86.6\,k {\scriptsize[86.1\,k, 87.3\,k]} & 55.8\,k {\scriptsize[55.4\,k, 56.2\,k]} & 1.68 {\scriptsize[1.67, 1.69]} & 1.09 {\scriptsize[1.08, 1.10]}\\
$2^{15}$ & 139.6\,k {\scriptsize[138.2\,k, 140.2\,k]} & 197.6\,k {\scriptsize[196.2\,k, 199.4\,k]} & 131.0\,k {\scriptsize[128.1\,k, 133.0\,k]} & 1.42 {\scriptsize[1.41, 1.44]} & 0.95 {\scriptsize[0.92, 0.97]}\\
$2^{16}$ & 400.2\,k {\scriptsize[397.2\,k, 407.1\,k]} & 526.2\,k {\scriptsize[523.5\,k, 532.6\,k]} & 282.6\,k {\scriptsize[281.4\,k, 284.5\,k]} & 1.34 {\scriptsize[1.32, 1.35]} & 0.70 {\scriptsize[0.69, 0.71]}\\
$2^{17}$ & 1.000\,M {\scriptsize[992.2\,k, 1.008\,M]} & 1.858\,M {\scriptsize[1.832\,M, 1.871\,M]} & 701.9\,k {\scriptsize[698.5\,k, 713.0\,k]} & 1.85 {\scriptsize[1.83, 1.86]} & 0.71 {\scriptsize[0.70, 0.71]}\\
$2^{18}$ & 2.122\,M {\scriptsize[2.110\,M, 2.149\,M]} & 4.047\,M {\scriptsize[4.022\,M, 4.090\,M]} & 1.501\,M {\scriptsize[1.488\,M, 1.511\,M]} & 1.91 {\scriptsize[1.89, 1.93]} & 0.71 {\scriptsize[0.70, 0.71]}\\
$2^{19}$ & 4.491\,M {\scriptsize[4.360\,M, 4.675\,M]} & 8.844\,M {\scriptsize[8.725\,M, 8.946\,M]} & 5.195\,M {\scriptsize[5.158\,M, 5.242\,M]} & 1.97 {\scriptsize[1.94, 2.02]} & 1.21 {\scriptsize[1.17, 1.22]}\\
$2^{20}$ & 12.499\,M {\scriptsize[12.242\,M, 12.947\,M]} & 19.083\,M {\scriptsize[18.824\,M, 19.385\,M]} & 10.350\,M {\scriptsize[10.222\,M, 10.491\,M]} & 1.53 {\scriptsize[1.50, 1.56]} & 0.85 {\scriptsize[0.84, 0.87]}\\
\bottomrule\end{tabular}\end{table}

\begin{table}[h]\centering
\caption{Vector width, complex float: fastfft compiled for 128, 256, 512-bit vectors
(\texttt{FASTFFT\_SIMD\_BYTES}), median ns per transform [95\,\% CI]; SimdFFT for reference.}\label{tab:width}
\begin{tabular}{rllll}\toprule
$N$ & 128 bit & 256 bit & 512 bit & SimdFFT\\\midrule
$2^{6}$ & 64.4 {\scriptsize[64.1, 66.4]} & 50.0 {\scriptsize[49.5, 50.5]} & 35.0 {\scriptsize[34.7, 35.2]} & 114 {\scriptsize[112, 116]}\\
$2^{7}$ & 135 {\scriptsize[135, 136]} & 105 {\scriptsize[104, 106]} & 69.0 {\scriptsize[68.5, 69.7]} & 237 {\scriptsize[235, 238]}\\
$2^{8}$ & 333 {\scriptsize[331, 335]} & 221 {\scriptsize[220, 222]} & 162 {\scriptsize[161, 164]} & 468 {\scriptsize[464, 471]}\\
$2^{9}$ & 690 {\scriptsize[685, 696]} & 445 {\scriptsize[442, 449]} & 302 {\scriptsize[298, 305]} & 1018 {\scriptsize[1010, 1031]}\\
$2^{10}$ & 1502 {\scriptsize[1492, 1513]} & 980 {\scriptsize[977, 986]} & 655 {\scriptsize[649, 659]} & 2147 {\scriptsize[2134, 2159]}\\
$2^{11}$ & 3323 {\scriptsize[3313, 3340]} & 2332 {\scriptsize[2300, 2375]} & 1793 {\scriptsize[1784, 1807]} & 4942 {\scriptsize[4917, 4959]}\\
$2^{12}$ & 8781 {\scriptsize[8742, 8819]} & 6660 {\scriptsize[6515, 7907]} & 4895 {\scriptsize[4866, 4922]} & 12.2\,k {\scriptsize[12.1\,k, 12.2\,k]}\\
$2^{13}$ & 18.6\,k {\scriptsize[18.5\,k, 18.7\,k]} & 13.4\,k {\scriptsize[13.4\,k, 13.6\,k]} & 10.1\,k {\scriptsize[9997, 10.2\,k]} & 29.3\,k {\scriptsize[29.2\,k, 29.4\,k]}\\
$2^{14}$ & 40.6\,k {\scriptsize[39.3\,k, 40.8\,k]} & 28.2\,k {\scriptsize[27.9\,k, 28.9\,k]} & 20.5\,k {\scriptsize[20.3\,k, 20.6\,k]} & 59.1\,k {\scriptsize[58.8\,k, 59.5\,k]}\\
$2^{15}$ & 129.7\,k {\scriptsize[129.2\,k, 130.2\,k]} & 73.8\,k {\scriptsize[73.5\,k, 74.1\,k]} & 45.6\,k {\scriptsize[45.2\,k, 46.2\,k]} & 134.7\,k {\scriptsize[134.2\,k, 135.7\,k]}\\
$2^{16}$ & 296.3\,k {\scriptsize[294.1\,k, 299.4\,k]} & 201.6\,k {\scriptsize[200.1\,k, 203.7\,k]} & 140.9\,k {\scriptsize[140.4\,k, 141.7\,k]} & 271.7\,k {\scriptsize[269.7\,k, 273.6\,k]}\\
$2^{17}$ & 642.0\,k {\scriptsize[639.1\,k, 645.7\,k]} & 458.7\,k {\scriptsize[454.1\,k, 464.8\,k]} & 418.3\,k {\scriptsize[410.7\,k, 423.7\,k]} & 726.4\,k {\scriptsize[719.1\,k, 739.8\,k]}\\
$2^{18}$ & 1.432\,M {\scriptsize[1.423\,M, 1.448\,M]} & 1.083\,M {\scriptsize[1.072\,M, 1.089\,M]} & 1.008\,M {\scriptsize[1.001\,M, 1.015\,M]} & 1.808\,M {\scriptsize[1.792\,M, 1.832\,M]}\\
$2^{19}$ & 3.753\,M {\scriptsize[3.724\,M, 3.778\,M]} & 2.538\,M {\scriptsize[2.520\,M, 2.548\,M]} & 2.068\,M {\scriptsize[2.059\,M, 2.082\,M]} & 4.355\,M {\scriptsize[4.319\,M, 4.402\,M]}\\
$2^{20}$ & 8.402\,M {\scriptsize[8.341\,M, 8.485\,M]} & 5.956\,M {\scriptsize[5.845\,M, 5.999\,M]} & 5.524\,M {\scriptsize[5.439\,M, 5.616\,M]} & 9.040\,M {\scriptsize[8.948\,M, 9.148\,M]}\\
\bottomrule\end{tabular}\end{table}

\begin{table}[h]\centering
\caption{Standard library: complex float built with GCC/libstdc++ (native simd) and with clang/libc++ (vir-simd array
emulation), median ns per transform [95\,\% CI].}\label{tab:libcxx}
\begin{tabular}{rllll}\toprule
$N$ & fastfft, libstdc++ & fastfft, libc++ & SimdFFT, libstdc++ & SimdFFT, libc++\\\midrule
$2^{6}$ & 35.0 {\scriptsize[34.7, 35.2]} & 162 {\scriptsize[160, 162]} & 114 {\scriptsize[112, 116]} & 226 {\scriptsize[225, 227]}\\
$2^{7}$ & 69.0 {\scriptsize[68.5, 69.7]} & 721 {\scriptsize[716, 732]} & 237 {\scriptsize[235, 238]} & 461 {\scriptsize[457, 465]}\\
$2^{8}$ & 162 {\scriptsize[161, 164]} & 1382 {\scriptsize[1375, 1392]} & 468 {\scriptsize[464, 471]} & 879 {\scriptsize[871, 887]}\\
$2^{9}$ & 302 {\scriptsize[298, 305]} & 2763 {\scriptsize[2752, 2771]} & 1018 {\scriptsize[1010, 1031]} & 1961 {\scriptsize[1954, 1977]}\\
$2^{10}$ & 655 {\scriptsize[649, 659]} & 6141 {\scriptsize[6097, 6186]} & 2147 {\scriptsize[2134, 2159]} & 4039 {\scriptsize[4019, 4097]}\\
$2^{11}$ & 1793 {\scriptsize[1784, 1807]} & 18.0\,k {\scriptsize[17.7\,k, 18.4\,k]} & 4942 {\scriptsize[4917, 4959]} & 9468 {\scriptsize[9323, 9630]}\\
$2^{12}$ & 4895 {\scriptsize[4866, 4922]} & 36.1\,k {\scriptsize[35.9\,k, 36.2\,k]} & 12.2\,k {\scriptsize[12.1\,k, 12.2\,k]} & 20.8\,k {\scriptsize[20.7\,k, 20.9\,k]}\\
$2^{13}$ & 10.1\,k {\scriptsize[9997, 10.2\,k]} & 73.0\,k {\scriptsize[72.5\,k, 73.9\,k]} & 29.3\,k {\scriptsize[29.2\,k, 29.4\,k]} & 46.8\,k {\scriptsize[46.6\,k, 47.5\,k]}\\
$2^{14}$ & 20.5\,k {\scriptsize[20.3\,k, 20.6\,k]} & 158.3\,k {\scriptsize[157.5\,k, 160.1\,k]} & 59.1\,k {\scriptsize[58.8\,k, 59.5\,k]} & 97.8\,k {\scriptsize[96.2\,k, 99.2\,k]}\\
$2^{15}$ & 45.6\,k {\scriptsize[45.2\,k, 46.2\,k]} & 403.3\,k {\scriptsize[399.4\,k, 408.8\,k]} & 134.7\,k {\scriptsize[134.2\,k, 135.7\,k]} & 211.8\,k {\scriptsize[210.6\,k, 213.7\,k]}\\
$2^{16}$ & 140.9\,k {\scriptsize[140.4\,k, 141.7\,k]} & 839.8\,k {\scriptsize[831.3\,k, 853.3\,k]} & 271.7\,k {\scriptsize[269.7\,k, 273.6\,k]} & 444.1\,k {\scriptsize[441.4\,k, 452.7\,k]}\\
$2^{17}$ & 418.3\,k {\scriptsize[410.7\,k, 423.7\,k]} & 1.738\,M {\scriptsize[1.724\,M, 1.751\,M]} & 726.4\,k {\scriptsize[719.1\,k, 739.8\,k]} & 1.046\,M {\scriptsize[1.036\,M, 1.060\,M]}\\
$2^{18}$ & 1.008\,M {\scriptsize[1.001\,M, 1.015\,M]} & 3.664\,M {\scriptsize[3.646\,M, 3.685\,M]} & 1.808\,M {\scriptsize[1.792\,M, 1.832\,M]} & 2.322\,M {\scriptsize[2.307\,M, 2.348\,M]}\\
$2^{19}$ & 2.068\,M {\scriptsize[2.059\,M, 2.082\,M]} & 8.803\,M {\scriptsize[8.709\,M, 8.907\,M]} & 4.355\,M {\scriptsize[4.319\,M, 4.402\,M]} & 5.370\,M {\scriptsize[5.314\,M, 5.450\,M]}\\
$2^{20}$ & 5.524\,M {\scriptsize[5.439\,M, 5.616\,M]} & 19.276\,M {\scriptsize[19.058\,M, 19.409\,M]} & 9.040\,M {\scriptsize[8.948\,M, 9.148\,M]} & 11.230\,M {\scriptsize[11.106\,M, 11.333\,M]}\\
\bottomrule\end{tabular}\end{table}

\begin{table}[h]\centering
\caption{Accuracy: largest $\max_k|X_k-\hat X_k|/\max_k|\hat X_k|$ over all $N$, against a long-double reference.}\label{tab:acc}
\begin{tabular}{lccc}\toprule & fastfft & SimdFFT & FFTW\\\midrule
float & \num{2.1e-7} & \num{4.8e-7} & \num{2.4e-7}\\
double & \num{6.0e-16} & \num{8.4e-16} & \num{3.3e-16}\\
\bottomrule\end{tabular}\end{table}

\begin{table}[h]\centering
\caption{gnuradio4 blocks before and after the patch, gnuradio4's benchmark facility: median call time over runs,
speedup as median of paired per-run ratios [95\,\% CI].}\label{tab:gr4}
\begin{tabular}{llll}\toprule benchmark row & unpatched & patched & speedup\\\midrule
\texttt{std::complex$<$float$>$ - fft} & 70.5 \,\textmu s & 68.9 \,\textmu s & 1.06 {\scriptsize[0.98, 1.07]}\\
\texttt{std::complex$<$float$>$ - fftCT} & 79.7 \,\textmu s & 80.6 \,\textmu s & 1.00 {\scriptsize[0.95, 1.01]}\\
\texttt{std::complex$<$double$>$ - fft} & 93.0 \,\textmu s & 60.2 \,\textmu s & 1.55 {\scriptsize[1.45, 1.59]}\\
\texttt{std::complex$<$double$>$ - fftCT} & 98.0 \,\textmu s & 97.9 \,\textmu s & 0.99 {\scriptsize[0.91, 1.21]}\\
\texttt{float - fft} & 23.3 \,\textmu s & 22.7 \,\textmu s & 1.00 {\scriptsize[0.97, 1.07]}\\
\texttt{double - fft} & 27.4 \,\textmu s & 27.7 \,\textmu s & 1.01 {\scriptsize[0.81, 1.10]}\\
\texttt{FFT$<$float32$>$ stream N=     64} & 138.3 \,\textmu s & 53.5 \,\textmu s & 2.55 {\scriptsize[2.53, 4.34]}\\
\texttt{FFT$<$float32$>$ stream N=    256} & 135.3 \,\textmu s & 56.6 \,\textmu s & 2.39 {\scriptsize[2.33, 4.00]}\\
\texttt{FFT$<$float32$>$ stream N=   1024} & 150.4 \,\textmu s & 59.0 \,\textmu s & 2.52 {\scriptsize[2.49, 2.57]}\\
\texttt{FFT$<$float32$>$ stream N=   4096} & 210.3 \,\textmu s & 97.6 \,\textmu s & 2.14 {\scriptsize[2.08, 2.18]}\\
\texttt{FFT$<$float32$>$ stream N=  16384} & 256.6 \,\textmu s & 103.9 \,\textmu s & 2.48 {\scriptsize[2.37, 2.60]}\\
\texttt{FFT$<$float32$>$ stream N=  65536} & 300.1 \,\textmu s & 178.0 \,\textmu s & 1.70 {\scriptsize[1.58, 1.73]}\\
\texttt{FFT$<$float32$>$ stream N= 262144} & 1979.0 \,\textmu s & 1159.0 \,\textmu s & 1.76 {\scriptsize[1.64, 1.83]}\\
\texttt{FFT$<$float32$>$ stream N=1048576} & 9300.0 \,\textmu s & 5582.0 \,\textmu s & 1.66 {\scriptsize[1.65, 1.73]}\\
\texttt{FFT$<$float64$>$ stream N=     64} & 195.0 \,\textmu s & 144.3 \,\textmu s & 1.34 {\scriptsize[1.19, 1.43]}\\
\texttt{FFT$<$float64$>$ stream N=    256} & 228.3 \,\textmu s & 161.8 \,\textmu s & 1.42 {\scriptsize[1.32, 1.47]}\\
\texttt{FFT$<$float64$>$ stream N=   1024} & 274.7 \,\textmu s & 205.5 \,\textmu s & 1.34 {\scriptsize[1.16, 1.37]}\\
\texttt{FFT$<$float64$>$ stream N=   4096} & 367.3 \,\textmu s & 276.4 \,\textmu s & 1.33 {\scriptsize[1.26, 1.41]}\\
\texttt{FFT$<$float64$>$ stream N=  16384} & 444.5 \,\textmu s & 315.7 \,\textmu s & 1.43 {\scriptsize[1.35, 1.60]}\\
\texttt{FFT$<$float64$>$ stream N=  65536} & 616.6 \,\textmu s & 490.2 \,\textmu s & 1.33 {\scriptsize[1.16, 1.55]}\\
\texttt{FFT$<$float64$>$ stream N= 262144} & 4268.0 \,\textmu s & 2326.0 \,\textmu s & 1.84 {\scriptsize[1.74, 1.87]}\\
\texttt{FFT$<$float64$>$ stream N=1048576} & 19490.0 \,\textmu s & 11470.0 \,\textmu s & 1.74 {\scriptsize[1.66, 1.76]}\\
\bottomrule\end{tabular}\end{table}

\end{document}
